\section{Mandate: is the problem real enough to need a Federal statute?}
\label{sec:mandate}

The first question is whether there is a problem only Congress can solve. A member
will not spend political capital on a bill that duplicates existing law or chases a
hypothetical harm, so the mandate has to be specific. H. R. 9510 rests on three gaps
in current law, set out in the bill's findings, that no amount of agency guidance has
closed \cite{Kawchak2026HR9510Bill}.

The first is the embodiment problem and the sequencing gap. Physical AI oncology
systems place autonomous and semi-autonomous robots in direct contact with patients,
and an error in the machine-generated code that directs a robot cannot be caught
after the fact the way a paperwork error can; by the time the action executes, the
patient has already been touched. No Federal law regulates the order of operations,
so nothing today requires that verification happen before the code is generated. The
second is the financial-data gap. Existing law already makes a clinical
investigator's financial arrangements reportable through the Open Payments system
\cite{openpayments}, but no law prices, records, or discloses the cost of verifying
robot-patient interaction code, so the one number a budget committee would want does
not exist. The third is the cost asymmetry: an embodied failure is priced in lives
and litigation, while the verification gate that prevents it is priced in compute,
storage, and a bounded amount of human review. The gate's cost is small, fixed, and
measurable; the failure it prevents is large and unbounded.

The human stakes are not abstract. The foundational estimate that tens of thousands
of Americans die each year from preventable medical error has anchored two decades
of safety legislation \cite{kohn2000toerr}, and later analysis placed medical error
among the leading causes of death in the United States \cite{makary2016medical}. A
statute that makes a specific, automatable safety check mandatory before an
autonomous robot acts is a direct response to that record, not a speculative one.
\autoref{tab:mandate} states each gap and what it leaves unaddressed.

\begin{psgfig}
\node[psgone] (e) at (0,0) {Embodiment\\problem};
\node[psgtwo] (s) at (4.0,0) {Sequencing\\gap};
\node[psgthree] (f) at (8.0,0) {Financial-data\\gap};
\node[psgact] (n) at (4.0,-1.7) {The mandate\\for a statute};
\draw[psgedge] (e) to[out=-90,in=160,looseness=0.5] (n.north west);
\draw[psgedge] (s)--(n.north);
\draw[psgedge] (f) to[out=-90,in=20,looseness=0.5] (n.north east);
\end{psgfig}
% \vspace{-0.05cm}
\psgcaption{\textbf{Figure 2.} Three gaps in existing law compose the mandate: the
embodiment and sequencing gap, the financial-data gap, and the cost asymmetry that
makes the gate cheap relative to the harm it prevents.}

\needspace{10\baselineskip}\begin{table}[H]
\footnotesize
\begin{tabularx}{\textwidth}{@{}L{3.4cm} Y@{}}
\toprule
\textbf{Gap in existing law} & \textbf{What it leaves unaddressed today}\\
\midrule
Embodiment and sequencing & An error in robot-patient interaction code cannot be caught in time once a robot is in contact with a patient, and no Federal law requires verification before the code is generated or executed.\\
Financial-data & No Federal law prices, records, or discloses the cost of verifying robot-patient interaction code, so no comparable number reaches a budget committee.\\
Cost asymmetry & Nothing in current law forces the small, fixed cost of a verification gate to be weighed against the large, unbounded cost of an embodied failure.\\
\bottomrule
\end{tabularx}
\caption{The mandate: the gaps in existing law that only a Federal statute closes.}
\label{tab:mandate}
\end{table}

The mandate question is answered not by alarm but by specificity: three named gaps,
each closed by a single new section, each tied to a harm a member can explain to a
constituent who lost someone they love to a preventable error.
